\section{Related Work}
\label{sec:related}


Synthesizable molecular design methods can be roughly divided into two categories. 
The first directly searches the combinatorial space of synthetic pathways \citep{vinkers2003synopsis,hartenfeller2012dogs,korovina2020chembo,gottipati2020learning,horwood2020reactor,bradshaw2020barking,gao2021amortized,swanson2024synthemol,cretu2024synflownet,seo2024generative,koziarski2024rgfn,zhu2025syngfn}. 
Their performance remains limited due to several unaddressed challenges.
First, the action space of building block selection is extremely large, the size of which often reaches hundreds of thousands or millions \citep{enamineBuildingBlocks}.
They generate synthetic pathways by explicitly selecting building blocks from such large libraries, which is difficult to sufficiently explore chemical space during both training and sampling.
Second, the search space of synthetic pathways is sparse as not all combinations of building blocks and reactions lead to valid syntheses.
Lastly, the representation of synthetic pathways is not informed of the structural features of the resulting product molecules.
Models that learn distributions over synthetic pathways are typically only aware of the combinatorial rules of reactions and building blocks.
As a result, the structural features of the products, which are crucial determinants of molecular properties, are not captured, leading to a gap between pathway generation and property optimization.
These issues limit the sampling efficiency of this category of methods, making them even less practical for real-world applications that require expensive property evaluations or involve multiple property constraints.

The second category trains models to construct synthetic pathways from input molecular graphs \citep{luo2024projecting,gao2025generative,sun2024procedural,sun2025synllama,lee2025rethinking}.
These models, however, cannot generate molecules conditioned on property specifications and remain limited in both chemical space coverage and efficiency.
Beyond these two categories, other methods focus on specific problem classes, such as structure-based drug design \citep{jocys2024synthformer,rekesh2025syncogen}, or formulate the task differently, for instance, treating synthesizability as an optimization objective \citep{guo2024saturn,guo2025directly}.
